\documentclass[11pt,a4paper]{article}
\usepackage[T1]{fontenc}
\usepackage[utf8]{inputenc}
\usepackage{lmodern}
\usepackage{amsmath,amssymb}
\usepackage[margin=25mm]{geometry}
\usepackage[hidelinks]{hyperref}
\hypersetup{pdftitle={Constructing Newton motion: refinement, continuation and regularity},pdfauthor={navstokgap research working notes}}
\newcommand{\tightlist}{\setlength{\itemsep}{0pt}\setlength{\parskip}{0pt}}
\setlength{\emergencystretch}{3em}
\setcounter{secnumdepth}{0}
\begin{document}
\hypertarget{constructing-newton-motion-refinement-continuation-and-regularity}{%
\section{Constructing Newton motion: refinement, continuation and
regularity}\label{constructing-newton-motion-refinement-continuation-and-regularity}}

\textbf{Maintained result.} For a globally Lipschitz acceleration,
arbitrary positive partitions and independently mixed kick--drift orders
construct the same unique initial-value motion, with an explicit error
proportional to the mesh. A conservative logarithmic force family shows
the exact limit of that sufficient hypothesis: an unbounded slope still
permits unique motion when \(0<\beta\le2\); finite-time arrival and
waiting branches occur for \(\beta>2\). Local action tends to zero in
both regimes. These are classical existence and uniqueness statements,
with no positive action premise.

The polygon proof uses standard stability, consistency and Gronwall
arguments. The arrival test is a direct energy calculation, not the
general first-order Osgood criterion. Teschl {[}@TeschlODE2012, §2.7,
pp.54--56 and Problem 2.25, p.58; passages read{]} supplies the Euler
compactness background; López Pouso {[}@LopezPouso2005; metadata and
abstract read{]} treats a broader second-order problem whose full
theorem is not used here. See the
\href{../docs/Teschl_ODE_2012.md}{source companion}. Internal Terra/Sol
checks of the written arguments were model reviews, not independent peer
review or novelty evidence.

\hypertarget{start-with-a-finite-newton-polygon}{%
\subsection{1. Start with a finite Newton
polygon}\label{start-with-a-finite-newton-polygon}}

Let \(m>0\) be the mass and \(f(q)=F(q)/m\) the acceleration. The
desired motion would satisfy
\[\dot q=v,\qquad\dot v=f(q),\qquad(q(0),v(0))=(q_0,v_0).\tag{1}\]

For now we do not assume that a solution exists. Given a positive time
step \(\tau\), define a drift and a force kick by \[
 D_\tau(q,v)=(q+\tau v,v),\qquad
 K_\tau(q,v)=(q,v+\tau f(q)).
 \tag{2}
\]

A step can kick first, \(D_\tau K_\tau\), or drift first,
\(K_\tau D_\tau\). The first choice uses the force at the old position;
the second uses the force at the new position. Both give a finite
piecewise straight trajectory with impulses. Neither is assumed to agree
exactly with the unknown continuous motion.

Fix \(T>0\), take any finite positive partition of \([0,T]\), choose
either ordering in each cell, and compose the maps chronologically.
There is no equal-step or nesting requirement. When a time is inserted,
recompute the later states from the fixed initial state and the
prescribed force rule.

Existence of the limit requires two kinds of estimates. Stability bounds
how differences grow. Consistency bounds the error in the finite force
balance. The two together will construct the ODE solution.

\hypertarget{lipschitz-control-supplies-stability}{%
\subsection{2. Lipschitz control supplies
stability}\label{lipschitz-control-supplies-stability}}

Assume for this construction that
\[|f(q)-f(Q)|\le\Lambda|q-Q|\quad(q,Q\in\mathbb R),\qquad\Lambda\ge0.
\tag{3}\]

The quantity \(\Lambda\) has units of inverse time squared. Choose a
reference frequency \(\rho>0\) and use the length-valued phase-space
norm \[
 \|(q,v)\|_\rho=|q|+\frac{|v|}{\rho},\qquad
 \alpha=\max(\rho,\Lambda/\rho),\quad
 \beta=\rho+\Lambda/\rho,\quad b=|f(0)|/\rho.
 \tag{4}
\]

There is no action constant in these definitions. Directly from (2), \[
 \operatorname{Lip}(D_\tau)\le1+\rho\tau,\qquad
 \operatorname{Lip}(K_\tau)\le1+\Lambda\tau/\rho.
 \tag{5}
\]

Thus either cell map has Lipschitz constant at most \(e^{\beta\tau}\).
The constants multiply under composition, and their exponents add. Every
finite polygon therefore has the same upper amplification
\(e^{\beta T}\), independently of the number and placement of vertices.

We also need a bound on the states themselves, because consistency
involves the force along the polygon. For either cell map \(M_\tau\), \[
 \|M_\tau z\|_\rho\le e^{\beta\tau}\|z\|_\rho
                         +\tau e^{\beta\tau}b.
 \tag{6}
\]

For a kick-first cell, the inhomogeneous term is actually at most
\(\tau(1+\rho\tau)b\); for a drift-first cell it is \(\tau b\).
Iterating (6) gives the common vertex bound
\[R=e^{\beta T}(\|z_0\|_\rho+Tb).\tag{7}\]

In particular \(|q_j|\le R\), \(|v_j|\le\rho R\), and
\(|f(q_j)|\le A:=|f(0)|+\Lambda R\). A very irregular partition has not
created a new force or stability constant.

\hypertarget{the-polygons-construct-the-ode-solution}{%
\subsection{3. The polygons construct the ODE
solution}\label{the-polygons-construct-the-ode-solution}}

Linearly interpolate both vertex coordinates to obtain
\(z_\pi(t)=(q_\pi(t),v_\pi(t))\). The displayed velocity is continuous;
it differs from the physical constant velocity along each straight
polygon segment by at most that cell's kick. This distinction disappears
with the mesh.

Let \(\delta=|\pi|\) be the maximum cell duration. In a kick-first cell,
with elapsed time \(s\in[0,\tau]\), \[
 \dot q_\pi-v_\pi=(\tau-s)f(q_j),\qquad
 \dot v_\pi-f(q_\pi)=f(q_j)-f(q_\pi).
 \tag{8}
\]

In a drift-first cell the analogous expressions are \(-sf(q_{j+1})\) and
\(f(q_{j+1})-f(q_\pi)\). The bounds in (7) give, for almost every time,
\[
 \|\dot z_\pi-\mathcal F(z_\pi)\|_\rho\le C\delta,
 \quad\mathcal F(q,v)=(v,f(q)),\quad C=|f(0)|+2\Lambda R.
 \tag{9}
\]

The vector field \(\mathcal F\) is Lipschitz with constant \(\alpha\) in
the norm (4). Two polygons with the same initial state consequently
satisfy \[
 \sup_{t\le T}\|z_\pi(t)-z_\sigma(t)\|_\rho
 \le C\frac{e^{\alpha T}-1}{\alpha}
                        (|\pi|+|\sigma|).
 \tag{10}
\]

To verify this estimate, integrate the two residual equations (9),
subtract, and apply the elementary integral Gronwall inequality. That
inequality follows by integrating the differential inequality for an
upper bound after multiplying by \(e^{-\alpha t}\).

Completeness of continuous path space now constructs a uniform limit
\(z(t)\). The residual integral tends to zero; the Lipschitz bound
passes \(\mathcal F(z_\pi)\) to \(\mathcal F(z)\) under the time
integral. Hence \[z(t)=z_0+\int_0^t\mathcal F(z(s))\,ds.\tag{11}\]

The limit is differentiable and solves (1). Applying the same difference
estimate with zero residual proves uniqueness. Letting one mesh tend to
zero in (10) gives a quantitative error for any other partition. Both
kick orders, mixed arbitrarily, lead to the same solution. Since \(T\)
was arbitrary, the global Lipschitz hypothesis constructs forward motion
for every finite time.

This is the first analogy with Balaban: one proves bounds for each
finite description before appealing to a limiting object. The limit
follows from estimates uniform in the refinement. Nothing in this
argument requires choosing \(\hbar>0\).

\hypertarget{local-continuation-and-velocity-dependent-forces}{%
\subsection{4. Local continuation and velocity-dependent
forces}\label{local-continuation-and-velocity-dependent-forces}}

For locally Lipschitz acceleration, extend it from a bounded region and
apply the construction until the trajectory leaves that region.
Continuation then requires bounds preventing escape from every compact
phase-space region.

For \(V\in C^1\) with locally Lipschitz \(\nabla V\), (1) has a unique
local solution, by the usual Cauchy--Lipschitz theorem. Along it,
\[E(q,v)=\tfrac12m|v|^2+V(q),\qquad \frac d{dt}E=0.\] If \(V\) is
bounded below, (2) bounds \(|v|\). On any finite interval
\(|q(t)|\le |q_0|+|t|\sup|v|\), so \((q,v)\) remains in a compact set.
The local continuation theorem then gives a unique solution for all
positive and negative times. Coercivity of \(V\) is a stronger condition
than is needed for this finite-time argument.

Smoothness alone does not guarantee global existence. For \(m=1\),
\(V(q)=-q^4/4\), \(q_0>0\) and \(v_0=q_0^2/\sqrt2\), the zero-energy
solution \[q(t)=\frac{q_0}{1-q_0t/\sqrt2}\] escapes to infinity at
\(t=\sqrt2/q_0\). A force singularity can also stop a classical
solution: for example, a radial attractive Kepler trajectory can reach
\(r=0\) in finite time. A rule for continuing through the singular point
is further data.

In three dimensions the Lorentz force, with a static scalar potential
and a static magnetic field, gives \[
 \dot q=v,\qquad m\dot v=-\nabla V(q)+e\,v\times B(q).
\] It is not a force depending on position alone. If \(\nabla V\) and
\(B\) are defined and locally Lipschitz on all of \(\mathbb R^3\), the
full phase-space vector field is locally Lipschitz, so local existence
and uniqueness still hold. Moreover \(v\cdot(v\times B)=0\), hence the
conservative energy is conserved. If \(V\) is bounded below, the same
speed and position bounds as above prove global existence and
uniqueness, even if \(B\) is not globally bounded.

For instance, if \(V\) is constant and \(B=B_0\hat z\), the field
rotates the transverse velocity with cyclotron frequency \(eB_0/m\) and
leaves the longitudinal velocity constant. A velocity-dependent force is
thus not an exception to the ODE theory; its dependence must be included
in the phase-space vector field being estimated.

If \(B=\nabla\times A\) on \(\mathbb R^3\), the quantum Hamiltonian is
\[H_{\hbar,A}=\frac1{2m}(-i\hbar\nabla-eA(q))^2+V(q).\] The kinetic
momenta generally do not commute:
\([\pi_i,\pi_j]=i\hbar e(\partial_iA_j-\partial_jA_i)\). A positive
scalar Brownian-bridge density does not directly cover this magnetic
problem: a magnetic Euclidean representation contains a phase. The
\href{quantum-blocks-construct-newton.md}{coherent-packet theorem}
treats the scalar Hamiltonian; it does not prove this magnetic
extension.

An endpoint problem fixes \(q(0)\) and \(q(T)\) rather than initial
velocity. It has a separate possible singularity. For the harmonic
oscillator the full phase-space derivative is \[
 \begin{pmatrix}
 \cos\omega T&\sin(\omega T)/\omega\\
 -\omega\sin(\omega T)&\cos\omega T
 \end{pmatrix},\qquad \det=1.
\]

The position response to initial velocity is \(\sin(\omega T)/\omega\).
It vanishes at a conjugate time while the full phase-space map stays
invertible. The corresponding endpoint action Hessian can fail to invert
even though the initial-value motion remains unique. The time-blocking
note gives a finite-polygon determinant that approaches this singularity
under refinement.

Therefore a failure of one position-only generating function is not, by
itself, a failure of Newton's dynamics. A physical completion argument
has to establish which extra information the complete theory needs and
why it cannot be retained classically.

\hypertarget{the-conservative-arrival-criterion}{%
\subsection{5. The conservative arrival
criterion}\label{the-conservative-arrival-criterion}}

Near zero let \(a\) be continuous, odd, positive on \((0,r_0)\), zero at
zero and smooth away from zero. Consider \[
 \ddot q=a(q),\qquad W(q)=\int_0^q a(s)\,ds\quad(q\ge0),
 \qquad V(q)=-mW(|q|),\quad m>0.
 \tag{L1}
\]

Assume a bounded continuous global extension of the force, smooth away
from zero, by cutting off the potential outside a neighborhood of zero.
The \href{newton-cusp-preparation-floor.md}{bounded-cusp argument},
section 1, then gives global existence and bounds preventing escape. The
issue remaining here is uniqueness, not existence.

Conservation of \(E=m\dot q^2/2+V(q)\) follows by the chain rule since
\(V\in C^1\). Any solution starting at rest at zero has \(E=0\). On a
positive departing segment its velocity must be \[
 \dot q=\sqrt{2W(q)},\qquad
 \mathcal T(q)=\int_0^q\frac{ds}{\sqrt{2W(s)}}.
 \tag{L2}
\]

\textbf{Arrival criterion.} If \(\mathcal T(q)=\infty\) for one, hence
every small \(q>0\), rest is the unique local solution from \((0,0)\).
If \(\mathcal T(q)<\infty\), the inverse \(Q=\mathcal T^{-1}\)
constructs departures \(q(t)=\pm Q(t-t_0)\) after every waiting time
\(t_0\).

For the first implication, a nonconstant solution has an excursion of
one sign starting from zero. Its velocity cannot vanish inside that
excursion because \(W>0\), and separation of variables would give a
finite value for the divergent integral (L2). For the converse,
\(Q'=\sqrt{2W(Q)}\) and \(Q''=a(Q)\). Both tend to zero at departure, so
adjoining rest gives a \(C^2\) Newton trajectory. Reversing the two
departures supplies two different incoming full states that merge at
zero in finite time, exactly as in the cusp note.

For \(a(q)=c q^\alpha\) this test reproduces the familiar \(\alpha<1\)
nonuniqueness threshold: its integrand is proportional to
\(q^{-(1+\alpha)/2}\). The logarithmic family below tests the more
delicate border where no power-law Lipschitz constant exists.

\hypertarget{exact-logarithmic-threshold}{%
\subsubsection{Exact logarithmic
threshold}\label{exact-logarithmic-threshold}}

Fix \(\Omega,\ell>0\), \(\beta>0\) and put \(r=\beta/2\). For
sufficiently small \(q>0\), with \(L=\log(\ell/q)>\max\{1,r\}\), set \[
 W(q)=\frac{\Omega^2}{2}q^2 L^\beta,\qquad
 a(q)=\Omega^2q\bigl(L^\beta-rL^{\beta-1}\bigr).
 \tag{L3}
\]

Extend \(a\) oddly near zero and the potential with a smooth cutoff
farther out. This gives a bounded \(C^1\) potential and continuous
bounded force; \(a(0)=0\) since \(q L^\beta\to0\). Its derivative on the
positive side is \[
 a'(q)=\Omega^2\left[L^\beta-\frac{3\beta}{2}L^{\beta-1}
              +\frac{\beta(\beta-1)}2L^{\beta-2}\right],
 \tag{L4}
\]

which diverges for every \(\beta>0\). Thus none of these forces is
locally Lipschitz at zero.

Nevertheless the energy-zero speed and arrival integral are exact: \[
 \dot q=\Omega qL^r,\qquad
 \mathcal T(q)=\frac1\Omega\int_L^\infty u^{-r}\,du.
 \tag{L5}
\]

Consequently rest is unique for \(\beta\le2\). For \(\beta>2\), \[
 \mathcal T(q)=\frac{L^{1-r}}{\Omega(r-1)},\qquad
 Q(t)=\ell\exp\left(-[\Omega(r-1)t]^{-1/(r-1)}\right).
 \tag{L6}
\]

The formula for \(Q\) is valid until the departure leaves the region
where (L3) is used. It is flat to every order at \(t=0\), so an exactly
resting body and one that leaves after a delay satisfy the same Newton
equation and initial data.

For \(\beta\le2\) the \emph{whole} extended IVP is unique. Away from
zero this is ordinary local Lipschitz uniqueness. At zero with nonzero
initial velocity, the conserved-energy first-order equation has a
nonzero square root and a \(C^1\) right-hand side near the crossing, so
it gives unique passage. At zero velocity the arrival criterion applies.
Bounded force then gives a global reversible point flow. Its dependence
on initial state is continuous: common finite-time bounds give
compactness for any converging initial sequence; every subsequential
limit solves the same unique IVP. Thus all subsequences have the same
limit. A divergent local slope has not destroyed reversible mechanics.

At the borderline \(\beta=2\), the outgoing energy-zero solution from
\(q_0>0\) is particularly informative: \[
 q(t)=\ell\left(\frac{q_0}{\ell}\right)^{e^{-\Omega t}},
 \qquad \dot q(0)=\Omega q_0\log(\ell/q_0).
 \tag{L7}
\]

While it remains in the logarithmic region, it stays continuous as
\(q_0\downarrow0\) at every fixed finite time, although its position
sensitivity diverges. This is a genuine example of unique stable motion
without a Lipschitz dependence bound at rest.

The general first-order Osgood sufficient test applied directly to a
modulus of size \(s[\log(1/s)]^\beta\) reaches only \(\beta\le1\). The
stronger range \(1<\beta\le2\) above uses the conservative Newton
structure and energy; failure of that first-order sufficient test is not
a proof of nonuniqueness.

\hypertarget{crossing-times-and-arbitrary-mesh-polygons}{%
\subsubsection{Crossing times and arbitrary-mesh
polygons}\label{crossing-times-and-arbitrary-mesh-polygons}}

Let \(q_n=\ell e^{-n}\) with \(n\) sufficiently large. Along the
zero-energy trajectory, the time for the annulus \(q_{n+1}<q<q_n\) is \[
 \tau_n=\frac1\Omega\int_n^{n+1}u^{-r}\,du,
 \qquad \tau_n\sim\frac1{\Omega n^r}.
 \tag{L8}
\]

On this annulus the local acceleration slope bound is of order
\(\Lambda_n=\Omega^2n^\beta\). The short-cell scale \(\Lambda_n^{-1/2}\)
is therefore of the same order as \(\tau_n\). One can subdivide an
annulus crossing into a fixed number of shorter steps to obtain a small
local Lipschitz amplification parameter. There is no uniform slope bound
across all annuli.

The issue is whether infinitely many controlled annulus crossings can
fit into finite physical time: \[
 \sum_{n\ge N}\tau_n
 \begin{cases}
 =\infty,&\beta\le2,\\
 <\infty,&\beta>2.
 \end{cases}
 \tag{L9}
\]

For \(\beta\le2\), a zero-energy trajectory cannot reach the
infinite-depth boundary at zero in finite time. For \(\beta>2\),
perfectly regular motion on each finite collection of annuli reaches
that boundary in finite time, and the branch datum is missing there.
Local estimates at every finite scale do not by themselves control their
infinite composition. This is the useful multiscale connection with the
Balaban lesson; it is not a transfer of his field-theory proof.

For the unique regime, arbitrary-mesh Euler polygons, and the lesson's
mixed kick--drift polygons, converge uniformly on every fixed finite
time interval. Bounded force gives common position and momentum bounds;
uniform continuity gives a vanishing integral residual; compactness
supplies solution subsequences; uniqueness makes the whole family
converge. No \(O(\text{mesh})\) rate is asserted here.

In the nonunique regime, exact zero initial data make these
deterministic polygons remain at rest on every partition. However,
preparations \(q_0\downarrow0\), \(\dot q_0=\Omega q_0L_0^r\downarrow0\)
have \[
 L(t)=\left[L_0^{1-r}+\Omega(r-1)t\right]^{-1/(r-1)}
       \longrightarrow [\Omega(r-1)t]^{-1/(r-1)}.
 \tag{L10}
\]

Thus they converge to the nonzero departure (L6) at fixed small \(t>0\).
For each positive \(q_0\), choose a fine enough polygon to approximate
that regular trajectory; a diagonal sequence then has vanishing mesh and
vanishing initial error but a departing limit. Refinement is not jointly
stable with preparation refinement at zero. This is a stronger failure
than an unbounded numerical error estimate, while keeping the exact-zero
polygon's resting limit explicit.

\hypertarget{vanishing-local-action-and-a-quantum-residual}{%
\subsubsection{Vanishing local action and a quantum
residual}\label{vanishing-local-action-and-a-quantum-residual}}

The positive outgoing momentum is \(p(q)=m\Omega qL^r\). Its canonical
action from zero to \(q\) is \[
 \mathcal J(q)=\int_0^q p(s)\,ds
       =m\Omega\ell^2\int_L^\infty e^{-2u}u^r\,du.
 \tag{L11}
\]

For \(L>r/2\), elementary bounds give \[
 \frac{m\Omega}2q^2L^r
 \le\mathcal J(q)
 \le\frac{m\Omega q^2L^r}{2-r/L}.
 \tag{L12}
\]

The lower bound uses \(u^r\ge L^r\); the upper bound uses
\((L+v)^r\le L^r e^{rv/L}\) after writing \(u=L+v\). Therefore
\(\mathcal J(q)\sim(m\Omega/2)q^2L^r\to0\) for every \(\beta>0\). The
temporal threshold (L9) creates no positive action threshold.

There is still a concrete quantum comparison. In the local region
\(V=-p(q)^2/(2m)\), take the stationary energy-zero test functions \[
 \psi_\pm(q)=p(q)^{-1/2}
          \exp\left(\pm\frac{i}{\hbar}\mathcal J(q)\right),
 \qquad H_\hbar=-\frac{\hbar^2}{2m}\frac{d^2}{dq^2}+V(q).
 \tag{L13}
\]

Their eikonal and first-derivative transport terms cancel exactly.
Writing \(b(q)=p(q)^{-1/2}\) gives
\(H_\hbar\psi_\pm =-(\hbar^2/2m)b''e^{\pm i\mathcal J/\hbar}\) and the
exact local ratio \[
 \frac{|H_\hbar\psi_\pm|}{|V\psi_\pm|}
 =\left[\frac{\hbar}{m\Omega q^2L^r}\right]^2
       \left|\frac34-\frac rL+\frac{r(r+2)}{4L^2}\right|
 \sim\frac{3\hbar^2}{16\mathcal J(q)^2}.
 \tag{L14}
\]

For example, the differentiation behind this identity is
\(p'/p=(1-r/L)/q\) and \(p''/p=-r[1-(r-1)/L]/(q^2L)\); then
\(b''/b=3(p'/p)^2/4-p''/(2p)\). Equation (L14) is a local residual
identity, not an error bound for a normalized state or the full
propagator.

At fixed \(\hbar>0\), action comparable with \(\hbar\) makes the local
residual of this trajectory ansatz order one even in the classically
unique regime. Every sufficiently small positive \(\hbar\) defines a
radius \(q_\hbar\) by \(\mathcal J(q_\hbar)=\hbar\), and that radius
tends to zero as \(\hbar\to0\). The bounded-potential Hamiltonian has
global unitary dynamics at each positive \(\hbar\) by the cusp note's
held Kato--Rellich input. These facts neither require \(\hbar=0\) to be
a member of the same state theory nor select a positive numerical lower
bound across the family. No relativistic ingredient or \(c\) is used.

\hypertarget{consequence-for-state}{%
\subsection{6. Consequence for STATE}\label{consequence-for-state}}

Initial-value construction, continuation, endpoint caustics and
finite-time mergers have different hypotheses. Unbounded force slope
alone establishes no merger and no action floor. The
\href{newton-cusp-preparation-floor.md}{cusp preparation theorem} uses
an actual full-state merger and additional preparation laws. The
separate \href{balaban-newton-time-blocking.md}{endpoint blocking
argument} retains inverse Hessians, normalization and source response;
the \href{newton-nonlinear-time-cell.md}{nonlinear cell} controls a
supplied positive weight. Neither operation follows solely from
constructing a classical orbit.
\end{document}
